\paragraph{C.3h (Known mechanism and proved exact graded gates on abelian varieties).}
Let $A$ be a $g$-dimensional abelian variety, $g\ge1$, over an algebraically
closed field, of arbitrary characteristic. All groups in this paragraph
have rational coefficients. Write
\[
 B^r(A)=\mathrm{CH}^r(A)_{\mathbb Q}/
                  \mathrm{CH}^r(A)_{\mathrm{alg},\mathbb Q},
 \qquad J^r(A)=\mathrm{CH}^r(A)_{\mathrm{alg},\mathbb Q}.
\]
Both inherit the Beauville grading, with index $s$ characterized by
$[m]^*x=m^{2r-s}x$. For an ample divisor class $h$, the following
are exact equivalences for this fixed abelian variety:
\begin{equation}\label{eq:C3-abelian-B}
 \begin{gathered}
 h^{g-2r}:B^r(A)\longrightarrow B^{g-r}(A)
 \text{ injective for every }2r\le g\\
 \Longleftrightarrow\quad B^*(A)_{(s)}=0\quad(s<0).
 \end{gathered}
\end{equation}
\begin{equation}\label{eq:C3-abelian-J}
 \begin{gathered}
 h^{g-2r+1}:J^r(A)\longrightarrow J^{g-r+1}(A)
 \text{ injective for every }2r\le g+1\\
 \Longleftrightarrow\quad J^*(A)_{(s)}=0\quad(s\le0).
 \end{gathered}
\end{equation}
The displayed implications do not assert these vanishing conditions
in general. The first is the adequate-equivalence version of
\citet[Proposition 5.11; footnote 1, p.~254]{rootBeauville2010};
the second applies the same representation to its invariant
algebraically trivial subgroup.

\begin{proof}
We give the string argument, including the converse directions.
Choose the symmetric ample class $\theta$ associated to $h$.
\citet[Theorem 4.2 and Proposition 5.2, pp.~260--261]{rootBeauville2010}
construct algebraic correspondences acting by
\[
 \begin{gathered}
 E(x)=\theta x,\qquad
 F(x)=\frac{\theta^{g-1}}{d(g-1)!}*x,\qquad
 d=\frac{\deg\theta^g}{g!},\\
 H(x)=(2r-g-s)x\quad(x\in\mathrm{CH}^r(A)_{(s)}).
 \end{gathered}
\]
where $*$ is the Pontryagin product. They form an
$\mathfrak{sl}_2$-representation that is a direct sum of finite
dimensional representations. All operators preserve algebraic
equivalence, since they are correspondences. Therefore they induce
the same kind of representation on $B^*(A)$ and on $J^*(A)$.
Each is a direct sum of strings
\[
        z,\theta z,\ldots,\theta^{g+s-2a}z,
 \qquad z\in M^a_{(s)},\quad Fz=0,\quad g+s-2a\ge0,
\]
where $M$ denotes either $B$ or $J$. The semisimplicity used here
is ordinary characteristic-zero representation theory over
$\mathbb Q$, independent of the characteristic of the base field.
It concerns these locally finite representations, not a
semisimplicity assertion about the cycle category.

In degree $r$, a nonzero term of this string is
$\theta^{r-a}z$ with $a\le r$ and $r-a\le g+s-2a$.
Multiplication by $\theta^k$ kills it exactly when
\[
                         r-a+k>g+s-2a.
\]
For $k=g-2r$, this inequality becomes $a>r+s$.
It is impossible when $s\ge0$. Conversely, if a negative
graded piece is nonzero, take a nonzero lowest element
$z\in B^a_{(s)}$ in one of its strings. Then
$2a\le g+s<g$ and $g-2a>g+s-2a$, so the asserted
map in degree $a$ kills $z$. This proves
\eqref{eq:C3-abelian-B} for $\theta$.

For $k=g-2r+1$, the corresponding inequality is
$a>r+s-1$, impossible when $s\ge1$. A nonzero string
with $s\le0$ has $2a\le g+s\le g$ and
$g-2a+1>g+s-2a$; its lowest element is therefore a
nonzero kernel element in an allowed degree. This proves
\eqref{eq:C3-abelian-J} for $\theta$.

Finally $h=T_b^*\theta$ for some $b\in A(k)$, as a rational
divisor class: their difference belongs to $\operatorname{Pic}^0(A)$,
and the polarization isogeny associated to $\theta$ is surjective
on points over the algebraically closed field. Clearing the
possible denominator in the symmetric class and using the
surjectivity of multiplication on $\operatorname{Pic}^0(A)(k)$
gives the rational-class assertion as well. The pullback
$T_b^*$ is an automorphism of both $B$ and $J$, and
\[
             h^k T_b^*x=T_b^*(\theta^k x).
\]
It conjugates the relevant maps, proving both claims for $h$.
On $B$ one may also use directly that translations act trivially
modulo algebraic equivalence.
\end{proof}

These formulas specify small remaining tests. For $g=5$, the only
possibly negative piece in degree two is $B^2(A)_{(-1)}$:
the Fourier correspondence identifies the $s$-piece there with
$B^{3+s}(A)_{(s)}$; for $s\le-2$ this is a divisor,
a codimension-zero group, or zero, and its negative part vanishes.
The degree-two $s=-1$ piece is entirely primitive, since its
lowering image is in the zero group $B^1(A)_{(-1)}$.
Hence the first nonidentity higher map in this dimension satisfies
\[
 \ker\bigl(h:B^2(A)\to B^3(A)\bigr)
                 \simeq B^2(A)_{(-1)}.
\]
For $g=4$ and $J^2$, the same Fourier argument eliminates
all negative pieces. Its degree-zero piece is primitive because
$J^1(A)_{(0)}=0$. Consequently
\[
 \ker\bigl(h:J^2(A)\to J^3(A)\bigr)
                 \simeq J^2(A)_{(0)}.
\]
For nonsymmetric $h$ the indicated identifications use a translation
as above and need not be canonical. Neither group on the right
is shown to vanish here. We have reduced these particular
injectivity problems to the stated graded groups, rather than
confusing them with the already refuted surjectivity assertions.
